	% Welcome! This is the unofficial University of Udine beamer template.

% We defined four theme colors: UniBrown, UniBlue, UniGold
% and UniOrange. For example, to write some uniud-brownish
% text, just use: \textcolor{UniBrown}{Hello!}

\documentclass[usenames,dvipsnames]{beamer}
\usepackage[utf8]{inputenc}
\usepackage{verbatim}
\usetheme{uniud}
%\bibliography{bib}
%%% Bibliography
%\usepackage[style=authoryear,backend=biber]{biblatex}
%\addbibresource{bib.bib}
\usepackage{tikz}
\usetikzlibrary{patterns,snakes}
%\DeclareNameAlias{author}{first-last}

%%% Suppress biblatex annoying warning
\usepackage{silence}
\WarningFilter{biblatex}{Patching footnotes failed}

%%% Some useful commands
% pdf-friendly newline in links
\newcommand{\pdfnewline}{\texorpdfstring{\newline}{ }} 
% Fill the vertical space in a slide (to put text at the bottom)
\newcommand{\framefill}{\vskip0pt plus 1filll}
\newcommand{\bigCI}{\mathrel{\text{\scalebox{1.07}{$\perp\mkern-10mu\perp$}}}}
\newcommand\asto{\mathrel{\stackrel{\makebox[0pt]{\mbox{\normalfont\tiny a.s.}}}{\to}}}


%ISMAEL'S COMMANDS
\newcommand{\given}{|}
\newcommand{\ra}{\rightarrow}
\newcommand{\rn}{\sqrt{n}}
\newcommand{\E}{\mathrm{E}}
\newcommand{\al}{\alpha}
\newcommand{\be}{\beta}
\renewcommand{\b}{\beta}
\newcommand{\ga}{\gamma}
\newcommand{\la}{\lambda}
\renewcommand{\l}{\lambda}
\newcommand{\epl}{\underline{\veps}_n}
\newcommand{\h}{\eta}
\renewcommand{\k}{\kappa}
\newcommand{\s}{\sigma}
\newcommand{\te}{\beta}
\newcommand{\q}{\beta}
\newcommand{\m}{\mu}
\newcommand{\veps}{\varepsilon}
\newcommand{\pli}{+\infty}
\newcommand{\A}{\cA}
\newcommand{\cA}{{\mathcal{A}}}
\newcommand{\cS}{{\mathcal{S}}}

\newcommand{\EM}{}
\newcommand{\ep}{\varepsilon}
\newcommand{\Ga}{\Gamma}

\newcommand{\si}{\sigma}

\newcommand{\cN}{\mathcal{N}}

\newcommand{\cT}{\mathcal{T}}
\newcommand{\cV}{\mathcal{V}}
\newcommand{\cY}{\mathcal{Y}}
\newcommand{\mc}{\mathop{\mathrm{ mc}} }
\newcommand{\argmin}{\mathop{\arg\min}}


\newcommand{\RR}{\mathbb{R}}

\title[LASSO is not Fully Bayesian]{‌\vspace{0.5cm}\\
LASSO is not Fully Bayesian}
\author[Behrad Moniri]{
  Behrad Moniri
  \pdfnewline
  \texttt{bemoniri@ee.sharif.edu}
}
\institute{\small{EE Department, Sharif University of Technology}}
\date{}
\begin{document}

\begin{frame}
\titlepage
\end{frame}


\begin{frame}
\frametitle{LASSO is not fully Bayesian.}

\begin{itemize}
	
	\item Data Generation Model: $Y = X\beta_0 + \epsilon$ where $\epsilon \sim \mathcal{N}(0, I)$.
	
	\item The LASSO optimization problem:
	
	\begin{equation}
	\label{EqLASSO} \hat{\be}^{\mathrm{LASSO}}_\la= \argmin_{\be\in\RR^p}
	\bigl[ \|Y-X\be\| _2^2 + 2\la\|\be\|_1
	\bigr]
	\end{equation}
	\item The posterior mode for the prior $\be_i \sim \mathrm{Laplace}(\la)$. Hence it has a Bayesian flavor.
	
	\item \textbf{Frequentist Optimality}: Can attain the (near) minimax rate $O(s \log n)$ for the square Euclidean loss over $s$-sparse signals, if $\lambda \approx \sqrt{2\log(n)}$.		
\end{itemize}
\end{frame}



\begin{frame}
	\frametitle{Not really Bayesian!}
	\begin{itemize}
		\item We will show that the full posterior does not contract at the same speed as its mode!
		
		\item Useless for uncertainty quantification, the central idea of Bayesian inference.
		
		\pause 
		\item The good behavior of the LASSO estimators must be due to the  sparsity-inducing form of the posterior mode,\\ \textbf{not} the Bayesian connection. 
	\end{itemize}
	
\end{frame}

\begin{frame}
	\frametitle{The Main Result}
		\begin{alertblock}{Theorem (7) from [CSV15]} \label{lemlb}
			Assume that $X = I$. For any $\la=\la_n$ such that $\rn/\la_n\to
			\infty$, there exists $\delta>0$ such that, as $n\to\infty$,
			%
			\[
			\mathbb{E}_{\be^0=0} \Pi_{\la_n}^{\mathrm{LASSO}} \biggl(\be: \|\be
			\|_2\le\delta \sqrt{n} \biggl(\frac{1}{\la_n}\wedge1 \biggr) \Big| Y
			\biggr) \rightarrow0.
			\]
			%
		\end{alertblock}
	
		\begin{itemize}
			\item Let $\lambda_n = \sqrt{2\log(n)}$. The posterior places \textbf{no} weight on the ball $\|\be\|_2 = O\Big(\sqrt{\frac{n}{\log n}}\Big)$ which is substantially larger than the minimax rate $\sqrt{s \log(n)}$ {\tiny (unless the signal is extremely dense)}.
		\end{itemize}
\end{frame}


\begin{frame}
\frametitle{Lemma (5.2) form [CV12]}
We will first state two lemmas:
\begin{small}


\begin{exampleblock}{Lemma (5.2) form [CV12]} 
	For any prior probability distribution $\Pi_n$ on $\RR^n$,
	any positive measure $\tilde\Pi_n$ with $\tilde\Pi_n\le\Pi_n$,
	and any $\q_0$,
	%
	\[
	\int\frac{p_{\q}}{p_{\q_0}}(Y) \,d\Pi_n(\q) \ge\|\tilde\Pi_n\|
	e^{-\tilde
		\s^2/2+\tilde\m^T(Y-\q_0) },
	\]
	%
	where $\tilde\m=\int(\q-\q_0) \,d\tilde\Pi_n(\q)/\|\tilde\Pi_n\|$ and
	$\tilde\s^2=\int\|\q-\q_0\|^2_2 \,d\tilde\Pi_n(\q)/\|\tilde\Pi_n\|$.
	Also, for any $r>0$,
	%
	\[
	P_{\q_0} \biggl(\int\frac{p_{\te}}{p_{\te_0}}(Y) \,d\Pi_n(\te) \geq e^{-r^2}
	\Pi_n\bigl(\q:\|\q-\q_0\|_2<r \bigr) \biggr)\geq1 - e^{-r^2/8}.
	\]
	%
\end{exampleblock}
\end{small}
\end{frame}

\begin{frame}
\frametitle{Proof of Lemma (5.2): First Assertion}
\begin{small}
We can write:
\begin{align*}
\log\bigg(\int \frac{p_\beta}{p_{\beta_0}}(Y) \frac{d \Pi_n(\beta)}{\|\tilde \Pi_n\|}\bigg) &\geq \int \log\bigg(\frac{p_\beta}{p_{\beta_0}}(Y)\bigg) \frac{d \Pi_n(\beta)}{\|\tilde \Pi_n\|} && \text{(Jensen's)}\\
&\geq \int \log\bigg(\frac{p_\beta}{p_{\beta_0}}(Y)\bigg) \frac{d \tilde \Pi_n(\beta)}{\|\tilde \Pi_n\|}&& (\tilde\Pi_n \leq \Pi_n)\\
&= -\frac{\tilde\sigma^2}{2} + \tilde \mu^\top (Y-\b) && \text{(Gaussianity)}
\end{align*}
Where $\tilde\mu = \int (\b - \b_0) \frac{d\tilde\Pi_n}{\|\tilde\Pi_n\|}$ and  $\tilde\sigma^2 = \int \|\b - \b_0\|^2_2 \frac{d\tilde\Pi_n}{\|\tilde\Pi_n\|}$. 

Hence:

\[
\int\frac{p_{\q}}{p_{\q_0}}(Y) \,d\Pi_n(\q) \ge\|\tilde\Pi_n\|
e^{-\tilde
	\s^2/2+\tilde\m^T(Y-\q_0) }.
\]
\end{small}
\end{frame}

\begin{frame}
	\frametitle{Proof of Lemma (5.2): Second Assertion}
	\begin{small}
		\begin{itemize}

			\item
			Let $\tilde \Pi_n$ be  $\Pi_n$ restricted to the set $\{\b: \|\b-\b_0\|_2\leq r\}$. Trivially $\tilde\Pi_n \leq \Pi_n$, $\|\tilde\mu\| \leq r$, $\tilde\sigma \leq r$, and $\|\tilde \Pi_n\| = \Pi_n(\q:\|\q-\q_0\|_2<r)$.
			\item
			Under $P_{\b_0}$, we have that ${\tilde\mu}^\top (Y-\b_0)$ has the same distribution as  $\|\tilde \mu\|Z$ where $Z\sim\mathcal{N}(0, I_n)$. 
			\item
			Note that 
			\begin{equation*}
				P\bigg[Zr\leq -r^2 + r^2/2\bigg] \leq \exp\big(-r^2/8\big).
			\end{equation*}
			\item Hence:			
			\[
			P_{\q_0} \biggl(\int\frac{p_{\te}}{p_{\te_0}}(Y) \,d\Pi_n(\te) \geq e^{-r^2}
			\Pi_n\bigl(\q:\|\q-\q_0\|<r \bigr) \biggr)\geq1 - e^{-r^2/8}.
			\] 
			Which concludes the proof.\qed
		\end{itemize}
	\end{small}
\end{frame}


\begin{frame}
	\frametitle{Lemma (7.1) from [CV12]}
	
	\begin{exampleblock}{Lemma (7.1) from [CV12]}
		\label{lemlb}
		We have $\E_{\q_0}\Pi_n (\q:\|\q-\q_0\|<s_n\given Y
		)\to0$,
		for any $s_n$ for which there exist $r_n$ such that
		
		\[
		\frac{\Pi_n (\q:\|\q-\q_0\|< s_n )} {
			\Pi_n (\q:\|\q-\q_0\|< r_n )}=o\bigl(e^{-r_n^2}\bigr).
		\]
	\end{exampleblock}

\end{frame}

\begin{frame}
	\frametitle{Proof of Lemma (7.1)}
	\begin{small}
	Based on the previous lemma, there exists and event $\mathcal{A}_n$ with $\mathrm{Pr}[A_n]\geq 1-\exp(-r_n^2/8)$ such that \[
	\int\frac{p_{\te}}{p_{\te_0}}(Y) \,d\Pi_n(\te) \geq e^{-r^2}
	\Pi_n\bigl(\q:\|\q-\q_0\|<r \bigr).
	\]	
	\begin{align*}
		\E_{\q_0}\bigg[\Pi_n(\b: \|\b - \b_0\|<s_n|Y ) 1_{\mathcal{A}_n}\bigg] &= \E_{\q_0} \Bigg[\frac{\int_{\b: \|\b-\b_0\|<s_n}p_{\b}(Y)\; d\Pi_n(\b)}{\int p_\b (Y) d\Pi_n(\b)}  1_{\mathcal{A}_n} \Bigg]\\
		&\leq 
		\E_{\q_0} \Bigg[\frac{\int_{\b: \|\b-\b_0\|<s_n}\frac{p_{\b}}{p_{\b_0}}(Y)\; d\Pi_n(\b)}{e^{-r_n^2}\Pi_n(\b: \|\b-\b_0\|\leq r_n)}   \Bigg]\\
		&\leq \frac{\Pi_n(\b: \|\b-\b_0\|<s_n)}{e^{-r_n^2}\Pi_n(\b: \|\b-\b_0\|\leq r_n)}  \to 0.
	\end{align*}
	Which concludes the proof. \qed
		\end{small}
\end{frame}


\begin{frame}
	\frametitle{Proof of the Main Theorem}
	\begin{small}
	Assume that $\b_0 = 0$. We will use lemma (7.1) to prove this theorem. We will need to show that for  proper sequences $r_n$ and $s_n$, 
			\[
	\frac{\Pi_n (\q:\|\q\|_2< s_n )} {
		\Pi_n (\q:\|\q\|_2< r_n )}=o\bigl(e^{-r_n^2}\bigr).
	\]
	\vspace{-0.5cm}
	\begin{itemize}
		\item 
		Under the prior, $|\beta_i|\sim \mathrm{Exp}(\lambda)$ and $\|\beta\|_1\sim \mathrm{Gamma}(\lambda, n)$. Thus, 	
			\begin{align*}
			\Pi_n(\b: \|\b\|_2 < s_n) &\leq 	\Pi_n(\b: \|\b\|_1 < s_n\sqrt{n})\\
			&= \int_{0}^{\sqrt{n}s_n} \lambda^n u^{n-1} \frac{e^{-\lambda u}}{\Gamma(n)} du\leq \frac{(\lambda \sqrt{n} s_n)^n}{\Gamma(n+1)}.	
			\end{align*}	

		\item  		\vspace{-0.2cm}		 We can write,
			\begin{align*}
			\Pi_n(\b: \|\b\|_2 < r_n) &= \Big(\frac{\lambda}{2}\Big)^n \int_{\|b\|_2\leq r_n} e^{-\lambda \|b\|_1} \mathrm{d}b \geq \Big(\frac{\lambda}{2}\Big)^n e^{-\lambda \sqrt{n}r_n} v_n r_n^n,
			\end{align*}	
			where $v_n$ is the area of unit $l_2$ ball in $\RR^n$ and the last inequality follows from $\|\b\|_1 \leq \sqrt{n}\|\b\|_2$.
	\end{itemize}

	\end{small}
\end{frame}

\begin{frame}
	\begin{small}
	By substituting $v_n = \pi^{\frac{n}{2}}/\Gamma(1+\frac{n}{2})$, and last two inequalities, 
	
	\begin{align*}
		\frac{\Pi(\|\b\|_2<s_n)}{\Pi(\|\b\|_2<r_n)} &\leq \Big(\frac{2}{\sqrt{\pi}}\Big)^n \frac{\Gamma(\frac{n}{2}+1)}{\Gamma(n+1)} n^{n/2} \exp\Big(\lambda \sqrt{n}r_n - n\log\big(\frac{r_n}{s_n}\big)\Big)\\
		&\leq \exp\bigg(\lambda \sqrt{n}r_n - n\log(c\frac{r_n}{s_n})\bigg).
	\end{align*}
	
	Set $r_n = \sqrt{n}\big(\max(\lambda^{-1}, 1)\big)$ and $s_n = \delta r_n$. Under these assumptions, the conditions of lemma (7.1) hold and we have:
	$$\E_{\q_0}\Pi_n (\q:\|\q-\q_0\|<s_n\given Y,
	)\to0,$$
	equivalently,
	\[
	\mathbb{E}_{\be^0=0} \Pi_{\la_n} \biggl(\be: \|\be
	\|_2\le\delta \sqrt{n} \biggl(\frac{1}{\la_n}\wedge1 \biggr) \Big| Y
	\biggr) \rightarrow0.
	\]
	Which concludeds the proof of the main theorem. \qed
	\end{small}
\end{frame}


\begin{frame}
	\frametitle{References}
	\begin{footnotesize}
	\begin{enumerate}		
		\item [CSV15] Isma\"el Castillo, Johannes Schmidt-Hieber, and Aad van der Vaart, Bayesian Linear Regression with Sparse Priors, \textit{Annals of Statistics}, 2015.

		\item [CV12] Isma\"el Castillo, and Aad van der Vaart, \\		Needles and Straw in a Haystack: Posterior Concentration for Possibly Sparse Sequences, \textit{Annals of Statistics}, 2012.
	\end{enumerate}
	\end{footnotesize}
\end{frame}
\end{document}